%%=============================================================================
%% Conclusie
%%=============================================================================

\chapter{Conclusie}%
\label{ch:conclusie}

Deze bachelorproef onderzocht hoe Docker-containers die gelinkt zijn aan CTFd-challenges dynamisch beheerd, aangemaakt en beëindigd kunnen worden op basis van gebruikersacties en ingestelde parameters op on-premise infrastructuur. Op basis van de literatuurstudie, de verkennende standalone PoC en de uiteindelijke plugin-PoC kan geconcludeerd worden dat een plugin-gebaseerde aanpak binnen CTFd hiervoor een haalbare oplossingsrichting is.

%%----------------------------------------------------------------------
\section{Beantwoording van de onderzoeksvragen}%
\label{sec:conclusie-onderzoeksvragen}

De centrale onderzoeksvraag kan positief beantwoord worden. Het is mogelijk om vanuit CTFd challenge-containers dynamisch te starten en opnieuw te stoppen, zolang de containerlevenscyclus expliciet wordt gekoppeld aan challenge-acties en runtime-beleid. In de uitgewerkte PoC gebeurt dat via een eigen \texttt{containerized} challenge-type dat Docker-containers start met vooraf ingestelde limieten en ze achteraf ook weer opruimt.

Ook de deelvragen kunnen op basis van de resultaten beantwoord worden:

\begin{itemize}
  \item \textbf{Welke beveiligingsmaatregelen zijn minimaal vereist?} Een bruikbaar basisprofiel bestaat uit een read-only root filesystem, het droppen van Linux capabilities, \texttt{no-new-privileges}, CPU- en geheugenlimieten, een PID-limiet en netwerkisolatie per instantie.
  \item \textbf{Hoe worden resource-limieten en netwerkisolatie afgedwongen?} De PoC gebruikt Docker runtime-opties voor CPU, geheugen en PID-limieten en maakt per runtime-instantie een eigen bridge-netwerk aan.
  \item \textbf{Hoe communiceert CTFd met Docker?} In deze bachelorproef gebeurt dat via een eigen plugin in Python die de Docker SDK voor Python gebruikt om containers, netwerken en image-imports aan te sturen. Die communicatie werkt zowel tegen een lokale Docker-socket als tegen een remote Docker API op een andere host.
  \item \textbf{Hoe werd de stabiliteit gevalideerd?} Via een combinatie van geautomatiseerde tests, smoke tests, timeout-validatie, isolatie-experimenten, een beperkte loadtest met Locust en split-host validatie in zowel een lokale simulatie als een reële twee-VM-opstelling.
\end{itemize}

%%----------------------------------------------------------------------
\section{Beperkingen van de PoC}%
\label{sec:conclusie-beperkingen}

Hoewel de plugin-PoC de kern van de onderzoeksvraag beantwoordt, zijn er duidelijke beperkingen. De huidige implementatie werd wel gevalideerd in same-host modus en in een split deployment over twee VM's, maar ze toont nog niet hoe de oplossing zich gedraagt in een bredere meerhost-opstelling met taakverdeling over meerdere runtime-hosts. Ook reverse-proxy-routing, TLS-beveiligde remote Docker-communicatie, geavanceerde production hardening en aangepaste seccomp-profielen maken nog geen afgerond deel uit van de PoC.

Daarnaast was de standalone PoC nuttig als exploratief tussenresultaat, maar ze is niet bedoeld als eindarchitectuur. De standalone stap diende vooral om risico's vroeg te reduceren en technische keuzes te onderbouwen voordat de pluginintegratie werd gebouwd.

%%----------------------------------------------------------------------
\section{Aanbevelingen voor productiegebruik}%
\label{sec:conclusie-aanbevelingen}

Voor productiegebruik volstaat de huidige PoC nog niet zonder bijkomende maatregelen. Organisaties die verder willen bouwen op deze aanpak, moeten minstens rekening houden met:

\begin{itemize}
  \item hardening van de Docker-host en het netwerksegment waarin challenge-containers draaien;
  \item versleutelde en streng afgeschermde communicatie indien Docker niet lokaal op dezelfde host draait;
  \item een expliciet beheer van gepubliceerde poortbereiken en firewallregels op de runtime-host;
  \item bijkomende monitoring en alerting voor mislukte starts, cleanupfouten en capaciteitspieken;
  \item operationele procedures voor imagebeheer, updates en incidentrespons.
\end{itemize}

De plugin-aanpak blijft daarbij een passende basis, omdat ze de gebruikerservaring en het beheer in CTFd centraliseert en tegelijk voldoende controle biedt over runtime-beleid.

%%----------------------------------------------------------------------
\section{Mogelijke uitbreidingen}%
\label{sec:conclusie-uitbreidingen}

De belangrijkste volgende stap is het hardenen van de reeds bewezen split-host architectuur. Daarbij moeten minstens vier vragen beantwoord worden:

\begin{itemize}
  \item hoe remote Docker-communicatie via TLS en passende toegangscontrole concreet wordt opgezet;
  \item of een dunne tussenlaag toch wenselijk is tussen plugin en Docker-host voor extra beveiliging of foutafhandeling;
  \item wat de impact is van die opsplitsing op opstartlatency en cleanupbetrouwbaarheid onder hogere belasting;
  \item hoe het systeem reageert wanneer de Docker-host tijdelijk onbereikbaar is.
\end{itemize}

Andere interessante uitbreidingen zijn een reverse-proxy-gebaseerd toegangspunt, uitgebreidere loadtesting en verdere production hardening.
